\subsubsection{MCMF (Min Cost Max Flow)}
upd on 25.10.12 已严肃进行势能优化 (使用 SPFA 预处理 h[ ] 数组，后续使用 Dijkstra)。
\begin{lstlisting}
struct node{ll u,v,w,c,nx;}e[1000010];
ll hd[100010],cnt=1;
void add(ll u,ll v,ll w,ll c){e[++cnt]={u,v,w,c,hd[u]};hd[u]=cnt;}
void ad(ll u,ll v,ll w,ll c){add(u,v,w,c),add(v,u,0,-c);}
ll n,m,S,T,nodes;
ll cst;
void init_mcmf(int count,int source,int sink){
    nodes=count;S=source;T=sink;cnt=1;cst=0;fill(hd,hd+nodes+1,0);
}
ll pre[100010];
ll dis[100010];
ll h[100010];
const ll inf=LLONG_MAX/4;
ll arg(){
    ll u,v=T,fl=inf;
    while(v!=S)u=pre[v],fl=min(fl,e[u].w),v=e[u].u;
    v=T;while(v!=S)u=pre[v],e[u].w-=fl,e[u^1].w+=fl,v=e[u].u;
    cst+=fl*(dis[T]-(h[S]-h[T]));return fl;
}
struct myp{ll w,x;};
bool operator < (myp a,myp b){return a.w>b.w;}
bool in[100010];
bool dij(){
    priority_queue<myp>q;
    f(i,1,nodes)dis[i]=inf,pre[i]=in[i]=0;
    dis[S]=0;q.push({0,S});
    while(!q.empty()){
        ll u=q.top().x;q.pop();
        if(in[u])continue;in[u]=1;
        for(int i=hd[u];i;i=e[i].nx){
            if(e[i].w==0)continue;ll v=e[i].v;
            if(dis[u]+(e[i].c+h[u]-h[v])<dis[v]){
                dis[v]=dis[u]+(e[i].c+h[u]-h[v]);pre[v]=i;
                q.push({dis[v],v});
            }
        }
    }
    return dis[T]!=inf;
}
array<ll,2> mcmf(){
    if(S==T)return {0,0};
    ll fl=0;cst=0;
    while(dij()){
        fl+=arg();
        f(i,1,nodes)if(dis[i]!=inf)h[i]+=dis[i];
        // f(i,1,nodes)cout<<h[i]<<" ";cout<<endl;
    }return {fl,cst};
}
void spfa(ll s){
    queue<ll>q;
    fill(h,h+nodes+1,inf);
    fill(in,in+nodes+1,false);
    h[s]=0;q.push(s);in[s]=1;
    while(!q.empty()){
        ll u=q.front();q.pop();in[u]=0;
        for(int i=hd[u];i;i=e[i].nx){
            if(e[i].w==0)continue;ll v=e[i].v;
            if(h[u]+e[i].c<h[v]){
                h[v]=h[u]+e[i].c;pre[v]=i;
                if(!in[v])in[v]=1,q.push(v);
            }
        }
    }
}
int main(){
    n=d,m=d,S=d,T=d;init_mcmf(n,S,T);
    f(i,1,m){
        ll u=d,v=d,w=d,c=d;
        ad(u,v,w,c);
    }spfa(S);
    auto res=mcmf();
    cout<<res[0]<<" "<<res[1]<<endl;
    return 0;
}
\end{lstlisting}

\paragraph{使用条件与多测} 每组重建 hd、cnt=1、cst=0，读完全部边后 spfa(S) 再 mcmf()。仅更新有限距离点的势能，避免不可达点反复加 INF 溢出；固定残量网络增广不会使原来不可达的顶点突然可达。若另行增删边或换源，必须重算势能。初始源可达残量图不得含负费用环，否则 SPFA 可能不终止；任意负环网络须用后文负权边转化或消圈算法。容量、距离、势能、流量乘费用及总费用须满足 ll 范围，INF 必须大于所有合法最短路。复杂度为初始 SPFA 加 $O(A E\log V)$，A 是增广次数，非强多项式。
